\section{第三原则：共生原则}

沉淀区给出了位置，不完备性给出了理由，剩下的问题是这些位置与人之间是什么关系。共生容易被读成一种温和的道德表述，仿佛两种东西互相需要，于是应该彼此友好。这里的用法要窄一些，它描述的是一种能量形式上的互补：一方是可以在给定地形上高速移动的能力，另一方是必须靠时间积累才能形成的位置。第三原则处理的就是这种互补的分工，以及分工之后哪些阻力必须被保留。

\subsection{分工的本质：动能与势能}

动能取决于运动状态，势能取决于相互作用系统的位形。在近地面条件下，把1千克物体缓慢抬高1米，重力势能约增加9.8焦耳；完成得更快，需要在更短时间内提供相应功率。这里借用运动与位置的区别，考察可快速调用的能力与依赖长期关系形成的资格。后者的形成条件，需要从经验、承诺与制度中说明。

\subhead{可以瞬时供给的一侧}

信息处理也要支付物理代价。室温下擦除一个比特的兰道尔下限约为$3\times10^{-21}$焦耳，\cite{landauer1961} 与抬升重物所需的能量可以采用同一单位比较。差别在于操作对象、实际能耗和系统边界。复制模型、部署实例和执行动作分别需要计算、传输、存储与执行资源，某一环节变便宜，并不直接给出整条任务链的成本。

把这两个数字放回同一个比值里，量级差别就可以写成，

\begin{equation}
\frac{E_{\text{抬升}}}{E_{\text{擦除}}}=\frac{mgh}{k_{\mathrm{B}}T\ln 2}\approx 3\times10^{21}
\end{equation}

式中，$E_{\text{抬升}}=mgh$ 为把质量为 $m$ 的物体抬升高度 $h$ 所需的机械功，单位为焦；$m$ 为质量，单位为千克；$g$ 为重力加速度，单位为米每二次方秒；$h$ 为高度，单位为米；$E_{\text{擦除}}=k_{\mathrm{B}}T\ln 2$ 为擦除一个比特在热力学上允许的最小耗散，即兰道尔极限，单位为焦；$k_{\mathrm{B}}$ 为玻尔兹曼常数，单位为焦每开；$T$ 为环境温度，单位为开；$\ln 2$ 为 2 的自然对数。取正文所用的 1 千克、1 米与室温，右侧的数值约为 $3\times10^{21}$，也就是说，抬升一次重物所耗的能量相当于这个量级的比特擦除次数。这个式子的含义是：信息操作与物质做功不在同一个量级上，用算力总量去推断能改变多少物质，在物理上从一开始就不成立。适用条件需要一并给出：比值随质量、高度与环境温度变化，正文的 $3\times10^{21}$ 只对应它给定的那组取值；兰道尔极限是可逆擦除的理论下界，实际器件的能耗远高于它，因此这个比值给出的是量级判断，不是工程核算。

计算一侧的实际开销可以测量。Patterson等估计，训练GPT-3的电力消耗约为1287兆瓦时，相应碳排放估计约为552吨二氧化碳当量。\cite{patterson2021carbon} 这是特定训练任务与核算口径下的结果。训练、推理、实例部署和长期维护分别占用资源，能力的复制始终依赖供电、硬件与运行安排。

\subhead{必须积累的一侧}

再看必须积累的那一侧。人的位置大多不是由速度构成的，而是由一系列不可跳过的过程累积而成：技能的形成依赖重复到身体结构发生改变，判断力的校准依赖真实后果的反复反馈，信任的建立依赖多次守约，关系的内容依赖共同经历过的时间。这些过程的速率由生理和社会过程决定，不因为外部信息处理变快而变快。它们更接近势能：积累缓慢，一旦形成可以持续释放，而且不能在个体之间自由转移。这笔账也解释了为什么人无法通过“更快地学习”来追平前一节所述的那种速度优势：即使学习效率提高，能力形成的物理过程仍然需要时间，而这段时间本身就是势能的一部分。谁支付了这段时间，谁就获得了这段位置。

这类积累的速率不是估计值，行业标准把它写成了明文门槛。美国联邦航空管理局要求民航机长在取得执照前至少积累 1500 小时飞行时间，这条规则在 2009 年科尔根航空事故之后被写进法规，理由是低飞行时数的机组在突发情形下的处置能力不足。\cite{faa2013atp} 医疗行业的结构类似：一名外科医生从医学院毕业到独立主刀，通常需要五到七年的住院医师训练，专科方向还要更长。这些年限不能被购买，也不能被转让：执照可以考取，飞行小时数只能积累；资质可以被认证，手术的手感只能由自己的病例形成。势能的定价困难正来自这里，它没有替代品，因此无法通过增加供给来降价。

把两侧合起来，分工的内容就清楚了：一侧负责在给定地形上高速搜索，另一侧负责提供地形并承担地形错误带来的后果。这不是能力上的分工，因为定义地形并不需要更强的算力；它是位置上的分工，因为势能只能由承担了积累代价的一方持有。这一点和前一节的结论接得上：目标不能由系统内部生成，而外部的判据只能由那些已经支付了代价、并且会被自己判断错误所改变的位置提供。

\subhead{两种时间尺度的剪切应力}

两种时间尺度并置在一起，必然产生剪切应力。绘图的一侧以年为单位工作，登山的一侧以天甚至小时为单位工作，两者之间的速度差不会被技术消除，只会随着工具变快而扩大。它的表现形式都很具体：标准更新慢于能力变化，制度调整慢于市场节奏，个人的长期计划跟不上工具的换代周期，而在这些错位里，处在慢一侧的人会持续感到被拉扯。把这种张力说成需要消除的病灶是一种误读，因为消除它只有两条路，让慢的一侧加速或者让快的一侧停下：前者受内禀速率约束无法实现，后者要求放弃能力上的收益。可行的做法是承载张力的结构，把一部分判断留在慢的一侧，同时让快的一侧在明确的边界内运行。

\subhead{这条界线的限度}

这个分工必须接受几处限定，否则它会被读成一种身份宣告。动能与势能可以互相转化，所以这条界线不是天赋的界线：机器同样可以积累位置性资产，一座数据中心、一套被反复验证的流程、一个被市场接受的品牌，都具有势能的性质；人同样参与搜索，日常判断里大量成分是快速的模式匹配。真正的分界不在物种，而在代价能否转移：如果代价可以被转移、被保险、被外包、被回滚，那么位置就没有被真正承担，势能也就不成立。反向的限定同样成立：占据了一个位置不等于判断正确，势能提供的是资格而不是正确性。历史上由处在承受位置的人做出的灾难性判断并不少见，这个事实不能靠分工来抹掉。分工能够支持的只是一条程序上的要求：给出地形的人必须在这片地形上行走，并且能够被后果改变。

机器一侧的积累同样可以量化。一座先进制程晶圆厂的投资在百亿美元量级，建设周期以年计，投产后还要经过良率爬坡才能形成可以交付的产能；这套资产的重量体现在它不能快速迁移，也不能在需求变动时被随意改产。\cite{verheyde2019tsmc} 这些约束与算法进步无关，它们由建筑、设备、材料供应和工程人员的实际速度决定。位置性资产因此不是人类专属，它也不随模型迭代而瞬时重置：把一套能力复制到新的实例上只需要分钟，把一条产线建起来仍然需要数年。

\subsection{人作为绘图师：地形从哪里来}

问题已经落到目标本身从哪里来。要看清这个问题，最好从优化问题的写法入手。一个优化问题写作

\begin{equation}
x^{*} = \arg\max_{x \in \mathcal{X}} f(x)
\end{equation}

求解这个问题的意思是：在允许的集合 $\mathcal{X}$ 里，找出使目标函数 $f$ 取值最大的那个 $x$。算法工程的大部分工作是在这一层之下完成的：设计更好的搜索策略、更快的近似、更有效的表示，让 $x^{*}$ 被更快更准地找到。这个式子里有一项从来不参与求解，它只是被使用，那就是 $f$。目标函数不是优化过程产生的结果，而是优化过程的输入。一个算法能把 $f$ 最大化到小数点后多少位，是一个能力问题；$f$ 应该是什么，不属于这个问题的范围。

这个区分看起来是技术性的，但它决定了后面所有讨论的边界。当人们在讨论通用智能何时到来时，讨论的是求解能力那一侧；当人们在讨论这项技术会带来什么后果时，讨论的往往是 $f$ 那一侧，而那一侧并没有跟着一起进步。适应度景观是同一个结构在生物学中的表述：一个物种在景观上攀爬，向更高的适应度移动，但景观的形状由环境给出，不由攀爬者给出。\cite{kauffman1993} 机器优化路径，人定义地形。这个分工不是能力上的分工，因为定义地形并不需要比优化路径更强的算力，它需要的是另一种东西。

\subhead{地形从哪里来}

那么地形是什么，它从哪里来。地形不是一条可以直接读出的东西，它是一个社会在漫长时间里通过大量具体判断累积形成的。什么算进步，什么风险可接受，什么未来值得追求，什么代价不能转嫁，这些问题的答案从来没有被完整地写下来过，它们分布在法律、习俗、专业规范、公共舆论和一代人的共同经历里。这类答案有几个特征，使它们不能被自动化。

第一个特征是它们包含不可通约的比较。一项减少碳排放的技术和一项提高粮食产量的技术，哪一个更重要，这个问题没有一个共同单位可以度量。成本收益分析可以给出一个数字，但那个数字来自对两类价值的相对权重的预先设定，而权重本身就是争议所在，把它写成一个标量函数，只是把争议隐藏在了参数里。第二个特征是它们依赖代价的分布。同一个总收益，如果代价集中在少数无法议价的人身上，它和一个总收益略低但代价分散的方案的比较结果，取决于社会对公平的判断，这个判断不能由总效用来代替。这也是为什么基础设施选址、环境许可、药物审批这类问题无法被纯技术化处理：技术可以计算总量，不能决定分布。

第三个特征是它们随时间变化。三十年前被普遍接受的代价转嫁在今天不可接受，反过来也一样。地形的这一部分性质意味着目标函数不是一次给定就固定的，它需要在新的信息和新的承受经验中不断修订。这三个特征指向同一个结论：地形是社会在自己承受后果的过程中逐步给出的，它不是被发现的知识，而是被做出的决定，因此它带有决定的一切特征：有争议、会改变、需要有主体来承担。

\subhead{四类只能由承受者回答的问题}

把上面的一般论述落到具体，有四类问题在结构上必须由承受后果的主体来回答，而它们各自不能被技术评估替代的原因略有不同。什么算进步，这个问题的答案决定了资源的长期流向。以产出总量衡量进步和以人的能力发展衡量进步，会导向完全不同的技术路线。这个问题不能被交给一个只优化已有指标的机构，因为指标的设定本身就是这个问题的答案，把指标设定权交出去等于把问题本身交出去。

什么风险可以接受，这一类的困难在于阈值。每一项技术在带来收益的同时带来风险，风险的可接受程度取决于一个社会愿意让谁承担多少代价。核电站的选址、基因编辑的边界、自动驾驶的接管策略、新药的审批速度，这些问题都需要一个有权做出决定并且承担决定后果的主体；技术评估可以提供概率分布，不能提供接受阈值，因为从概率到一个“可以接受”的判断之间，跨过的是一段价值论证，不是一段计算。什么未来值得追求，这一类的困难在于时间尺度。载人航天、深空探测、基础科学、生态修复，这些方向在几十年尺度上都无法用投入产出比来辩护。选择一个长周期方向，本质上是让当代人为后代人的可能性支付成本，而这个选择必须由能够代表未来的人做出。

什么代价不能转嫁，是前面三个的底线形式。当一项技术把成本转移给没有参与决定的一方（转移给未来、转移给生态、转移给缺乏议价能力的群体），这个转嫁在技术上可能是最优的，在结构上是不可接受的。设定这条底线不需要技术知识，需要的是承认承受者有权拒绝，而“有权拒绝”这件事本身无法从任何优化过程里推导出来，它必须被作为一个前提写进去。

\subhead{这不是能力优先的立场}

决定目标涉及受影响者的资格。关于“由谁承受什么”的决定，需要让承担风险的一方参与，或由具有相应授权的机构代表。机器可以提供候选、预测和论证，参与者也需要借助专业知识；决定的合法性还须满足程序、权利与救济的要求。承担后果是一项重要条件，不能单独替代授权。

用这个标准看，绘图师的资格不来自计算能力、不来自知识总量、也不来自自我意识。它来自一个更朴素的事实：给出地形的人必须在这片地形上行走，他的判断会以他自己的代价被检验。一个系统如果在判断出错时不承担任何后果，它的判断就缺少校准来源，无论它的推理能力有多强。把 AI 放进目标设定的位置会有一个特殊困难。不是因为它算得不够好，而是因为它的输出和目标设定的正确性之间没有反馈回路。人类设定者判断错了，会在若干年后看到后果并以某种形式付出代价；一个没有承受边界的系统在这条链上是断开的，错误不会被回流，因而不会被修正。

需要补一句限定，否则这个立场会变得过强。承受位置给出的是资格，不是正确性。有人承受后果并且做出的判断仍然是错的，这是常态，历史上由承受者做出的灾难性判断数不胜数。这个论证能支持的只是程序上的要求：目标设定权应当归于承担后果的一方，并且这一方必须有能力被后果改变。它不支持任何关于判断质量的保证。

\subhead{绘图是个持续动作}

还有一个容易被误解的地方：绘图不是一次性给出一个目标，然后交给机器去执行。地形在真实过程中持续变化。技术改变了可行的选项集合，选项集合的改变又改变了价值排序，价值排序的改变使原本不可接受的代价变得可以接受，或者反过来。这个过程没有终点，也不会收敛到一个稳定的目标函数。扮演绘图师的角色，意味着持续参与这个过程，而不是在某一天宣布目标已经确定。

这件事在实践中表现为一系列不那么宏大的动作：参与标准制定，参与事故调查，参与伦理审查，参与资源分配的讨论，在技术评估中坚持把分布问题摆上桌面，在长期项目上保留说不的权力。这些动作大多缓慢、不产生可见成果，也无法通过更好的算法来加速。它们的作用不在于给出更好的答案，而在于保证答案始终有人负责。代价是清楚的：这意味着人不能把判断的责任交出去。当一项技术被证明比人做得更好时，把决定权交给它的诱惑是真实的，而且在很多局部问题上交出决定权确实是正确的做法。要区分的是这两类问题：在给定目标下的求解可以交出去，目标本身不行。这条线不是由技术能力划定的，它由承受结构划定。

承认绘图是持续动作之后，另一个更实际的问题立刻出现。既然阻力并不都值得保留，也不都应当删除，那么在具体的场景里，哪些阻力承担着筛选、校准或约束的功能，哪些只是历史留下的损耗。这个问题不能靠态度回答，它需要一个可以逐项执行的判据。

\subsection{必要摩擦：哪些阻力应当保留}

\subhead{为什么必须区分}

前面关于高黏滞位置承载价值的论证容易引出一个错误的结论：如果高黏滞位置承载价值，那么保留阻力就是好事。这个结论既不正确也没法操作，因为阻力从来不是一个整体，它是一堆性质完全不同的东西的集合。同一项技术可以在一次改动中同时消除三类完全不同的摩擦，而它们的后果差别极大。

设想一家医院把纸质病历换成电子系统。录入时间缩短了，这是效率改进；重复检查减少了，这是资源浪费的消除；但同时，医生在病历上写下的第一手判断被结构化字段替代，那些无法被选项覆盖的观察失去了记录位置，而事后复核时，审阅者看到的是符合规范的表单，不再是医生的原始思考。第三种变化和第二种在表面上都是“减少了工作量”，在结构上完全不同。如果把这三类变化都归为摩擦的减少，就会得出一个笼统的判断，既不能指导技术设计，也不能指导制度安排。所以需要一个分类，而这个分类要能落在一处一处具体的阻力上，而不是停在态度层面。

\subhead{三类摩擦}

第一类是无效摩擦。它由组织惯性、信息不对称、流程冗余、工具落后、沟通损耗产生，不承担任何检验功能。医院里同一份信息在五个系统里各录一遍，行政审批要跑三趟窗口，公司内部因为部门壁垒重复开发同一套工具，都属于这一类。这类摩擦没有存在的理由，消除它不产生任何结构性损失，只产生收益。需要注意的是，无效摩擦对身处其中的人往往并不显得无效：它有历史原因，有人靠它获得位置，也有真实的风险规避动机曾经支撑它，因此识别它需要追问它在今天还挡着什么，而不是它当初为什么被建起来。

第二类是工程摩擦。它由物理约束和技术成本产生，是真实存在的但可以被压缩或转移。混凝土需要养护时间，芯片制造需要洁净环境和多层光刻，跨洋运输需要燃料和时间。这类摩擦不能被消除，但可以通过更好的工艺、更优的调度、更合理的资源投入来降低。压缩它的代价通常以资源或时间的形式在其他地方出现：压缩养护时间要加早强剂或者改变配比，压缩物流时间要空运或者提高库存。工程摩擦的处理原则因此是显式的，比较压缩成本与收益，选择总代价更低的方案，这个比较可以算，也可以被审计。

第三类是必要摩擦。它由系统需要承担的功能产生，消除它会造成功能缺失，而这种缺失往往不会立刻显现。验证、承责、身体接触、长期关系、公共讨论、风险暴露、教育中的慢形成，都属于这一类。它们的共同特征是：表面上看起来像是效率损失，实际上承担着筛选、校准或约束的作用，一旦被移除，系统会在某个延迟之后失去某项能力。三类之中只有这一类无法在当期被量化，而这恰恰是它在管理决策中最容易被牺牲的原因。\key{成熟技术的标志不是消除一切阻力，而是有能力区分无效摩擦与必要摩擦。} 这个区分不是价值立场，它要求对每一处摩擦追问一个具体问题：它现在在筛掉什么，或者阻止什么。
\begin{figure}[ht]
\centering
% >>>>>>>>>> figs/fig-friction-classes.tex >>>>>>>>>>
% 图 6-1  三类摩擦：同一项改动会同时动到三种性质不同的阻力
% 三行的差别不在"多"与"少"，而在移除之后会发生什么。
% 左侧色块给出类别与后果显现的时延；右侧三栏依次是来源、处置与移除的后果。
% 判定的三个问题写在正文里，本图只承担分类本身：
%   问一 移除后系统会不会在别处重建等效机制；
%   问二 代价由谁承担、收益由谁获得；
%   问三 把它压缩到极限会发生什么。
\begin{tikzpicture}[x=1cm, y=1cm, font=\small,
  rw/.style={rounded corners=1.5pt, line width=0.6pt},
  ct/.style={font=\scriptsize, text=ink, align=center, anchor=center},
  hd/.style={font=\scriptsize\heiti, text=slate, anchor=south}
]
\useasboundingbox (-0.30,0.80) rectangle (13.25,-6.42);

\def\cl{0.15} \def\cA{2.55} \def\cB{7.05} \def\cC{10.55} \def\cD{12.95}
\def\mAX{1.35} \def\mBX{4.80} \def\mCX{8.80} \def\mDX{11.75}

% ---------- 表头 ----------
\node[hd] at (\mAX,0.36) {类别};
\node[hd] at (\mBX,0.36) {来源};
\node[hd] at (\mCX,0.36) {处置};
\node[hd] at (\mDX,0.36) {移除的后果};

% ---------- 第一行：无效摩擦 ----------
\fill[rust!13] (\cl,0.00) rectangle (\cD,-1.75);
\fill[rust!30] (\cl,0.00) rectangle (\cA,-1.75);
\draw[rw, rust!55] (\cl,0.00) rectangle (\cD,-1.75);
\draw[slate!30, line width=0.4pt] (\cA,0.00) -- (\cA,-1.75);
\draw[slate!30, line width=0.4pt] (\cB,0.00) -- (\cB,-1.75);
\draw[slate!30, line width=0.4pt] (\cC,0.00) -- (\cC,-1.75);
\node[font=\small\heiti, text=rust!85, anchor=center] at (\mAX,-0.70) {无效摩擦};
\node[font=\scriptsize, text=slate, anchor=center] at (\mAX,-1.22) {后果时延：当期};
\node[ct] at (\mBX,-0.875) {组织惯性、信息不对称、\\流程冗余、工具落后};
\node[ct] at (\mCX,-0.875) {直接删除};
\node[ct] at (\mDX,-0.875) {无结构性损失};

% ---------- 第二行：工程摩擦 ----------
\fill[terra!13] (\cl,-1.95) rectangle (\cD,-3.70);
\fill[terra!30] (\cl,-1.95) rectangle (\cA,-3.70);
\draw[rw, terra!60] (\cl,-1.95) rectangle (\cD,-3.70);
\draw[slate!30, line width=0.4pt] (\cA,-1.95) -- (\cA,-3.70);
\draw[slate!30, line width=0.4pt] (\cB,-1.95) -- (\cB,-3.70);
\draw[slate!30, line width=0.4pt] (\cC,-1.95) -- (\cC,-3.70);
\node[font=\small\heiti, text=terra!90, anchor=center] at (\mAX,-2.65) {工程摩擦};
\node[font=\scriptsize, text=slate, anchor=center] at (\mAX,-3.17) {后果时延：数月};
\node[ct] at (\mBX,-2.825) {物理约束与技术成本\\（养护、光刻、运输）};
\node[ct] at (\mCX,-2.825) {比较压缩成本与收益\\之后压缩};
\node[ct] at (\mDX,-2.825) {代价以资源或\\时间形式转移};

% ---------- 第三行：必要摩擦 ----------
\fill[forest!13] (\cl,-3.90) rectangle (\cD,-5.65);
\fill[forest!30] (\cl,-3.90) rectangle (\cA,-5.65);
\draw[rw, forest!60] (\cl,-3.90) rectangle (\cD,-5.65);
\draw[slate!30, line width=0.4pt] (\cA,-3.90) -- (\cA,-5.65);
\draw[slate!30, line width=0.4pt] (\cB,-3.90) -- (\cB,-5.65);
\draw[slate!30, line width=0.4pt] (\cC,-3.90) -- (\cC,-5.65);
\node[font=\small\heiti, text=forest!85, anchor=center] at (\mAX,-4.60) {必要摩擦};
\node[font=\scriptsize, text=slate, anchor=center] at (\mAX,-5.12) {时延：跨管理周期};
\node[ct] at (\mBX,-4.775) {系统必须承担的功能\\（验证、承责、公共讨论）};
\node[ct] at (\mCX,-4.775) {保留；删除须先给出\\替代机制};
\node[ct] at (\mDX,-4.775) {功能缺失，\\延迟显现};

% ---------- 底部：两类误判的不对称 ----------
\node[font=\scriptsize, text=ink, anchor=north west, align=left]
      at ({\cl},{-6.02})
      {两个方向的误判不对称：保留无效摩擦是慢性损耗，删除必要摩擦是结构损伤。};
\end{tikzpicture}

% <<<<<<<<<< figs/fig-friction-classes.tex <<<<<<<<<<

\caption{三类摩擦。分类的依据不是量的多少，而是移除之后会发生什么；右端一栏给出三种后果，它们显现的时延自上而下依次拉长。这解释了为什么必要摩擦在当期管理决策中最容易被牺牲：它的代价要跨过管理周期才浮现，而收益的账面在当期。}
\label{fig:friction-classes}
\end{figure}

\subhead{判定的三个问题}

这个追问可以拆成三个可以逐项回答的问题。第一个问题是：如果移除它，系统会不会在别处重建一个等效机制。如果会，它属于无效摩擦或工程摩擦，因为系统并不真正依赖它。航空业的行李安检曾经由人工开箱完成，后来改成 X 光机加抽查，效率提高了几倍，安全水平没有下降，因为检查的功能被更高效的手段承担了。相反，学术界的同行评议如果被移除，替代机制必须提供同样的东西：让工作在发表前经过有资格的反对者检验。目前没有更便宜的替代方案，所以它尽管慢、尽管有偏见，仍然属于必要摩擦。

第二个问题是：它的代价由谁承担，收益由谁获得。这一类摩擦如果承担者和受益者不是同一方，通常意味着它是被转嫁的成本，而不是被设计的功能。平台把内容审核的成本转移给创作者和用户，用户的注意力成本被转移给广告主，这类结构中的“低摩擦”往往只是把摩擦挪到了看不见的地方。反过来，如果承担者和受益者是同一方，而且代价换来了真实的信息或能力，它更可能是必要的。这个问题的价值在于它能把讨论从“效率高低”转到“代价落在谁身上”，而后者才是真正需要被裁决的部分。

第三个问题是：把它压缩到极限会发生什么。工程摩擦在这个测试下表现良好，压缩到极限无非是成本上升，不会改变系统的性质。必要摩擦在这个测试下会失效：验证压缩到零，系统就不再知道什么是真的；教育中的慢形成压缩到零，得到的是通过考试的人，不是具备能力的人；公共讨论压缩到零，得到的是快速共识，不是被检验过的共识。这个问题之所以好用，是因为它不需要知道摩擦的成因，只需要设想一个极限情形。三个问题不能完全机械化，但比“阻力越大越好”或者“越快越好”都要可操作。

\subhead{与反技术立场的分歧}

把必要摩擦单列出来，容易被读成一种反技术的姿态，因此需要把分歧说清楚。Theodore Kaczynski 在 1995 年提出过一个论证：工业技术系统会持续扩张，逐步剥夺个人自主，因此应当被瓦解。\cite{kaczynski1995industrial} 他的论证中有一部分值得认真对待：技术系统确实会重新组织人的生活，而且这种重组往往不经过个人同意。但他的结论把技术当作一个整体来处理，因此只能得出全盘接受或全盘拒绝两个选项，而这两个选项在实践中都不可用。

必要摩擦的分析给出的是一条不同的路径。它不追问技术应不应该存在，而追问每一处具体的阻力承担了什么功能。这个问法允许同时做出两个方向相反的判断：一套审批流程应当被彻底删掉，而同一天上线的另一项自动化应当被紧急叫停。它不承认“技术本身”这个主语，只承认具体的机制在具体的位置上做了什么事。这个立场的代价是它不提供简洁的口号。它不能给出一个总的方向，只能给出一个逐项判断的程序，而逐项判断需要领域知识、需要理解系统的实际功能、需要在信息不完整的情况下做出取舍。这使它在公共讨论中处于劣势：反对技术的人有清晰的态度，支持技术的人有清晰的收益，管理摩擦的人只有一份长长的、需要专业判断的清单。

\subhead{两种误分类}

实践中最常见的错误不是分不清三类摩擦，而是在两个方向上系统性地误判，而这两个方向的错误并不对称。第一个方向是把必要摩擦当成无效摩擦删掉。这类错误几乎总是以效率为名，而且短期效果良好。企业取消资深员工的复核环节以提高交付速度，头几个月指标上升，一年后事故率上升，而事故的成因很难追溯到那次流程简化。教育压缩实践环节以增加课时覆盖，学生的考试成绩可能提高，动手能力在毕业后才暴露问题。这类错误的共同结构是：功能被移除，损失延迟出现，因果链在时间上被拉长到超出管理周期。

第二个方向是把无效摩擦当成必要摩擦保留。这类错误的辩护词通常是“传统”“经验”“风险”，它在稳定的组织里特别常见，因为保留摩擦对在位者有利。这类错误的成本是持续的、可见的、容易被抱怨的，所以相对容易被纠正，只要有足够的外力推动。两个方向的错误因此不对称：保留无效摩擦是慢性损耗，删除必要摩擦是结构损伤。这个不对称意味着，在信息不足的情况下默认选择应当是谨慎的，但这不等于把所有阻力都保留。谨慎的正确含义是：对每一处待删除的摩擦，要求提出者说明它在筛掉什么；如果提不出，就删掉；如果提得出，就要求给出替代机制。

到这里，一个自然的问题出现了。既然摩擦需要被逐项区分，而区分的判据又要求承担者与受益者对齐，那么谁来完成这件事。个体面对的是一个持续把人推向最短路径的环境，组织面对的是当期指标与长期能力之间的取舍。这两层都不是分类问题，而是实践问题。
